% Generated draft: requires mathematical and source review.
\documentclass[11pt]{article}
\usepackage[margin=25mm]{geometry}
\usepackage{fontspec}
\setmainfont{Latin Modern Roman}
\usepackage{amsmath,amssymb,mathtools,unicode-math}
\setmathfont{Latin Modern Math}
\usepackage{enumitem,needspace,xurl,hyperref}
\setlength{\parindent}{0pt}
\setlength{\parskip}{5pt}
\setlength{\emergencystretch}{3em}
\newcommand{\N}{\mathbb{N}}
\newcommand{\Z}{\mathbb{Z}}
\newcommand{\Q}{\mathbb{Q}}
\newcommand{\R}{\mathbb{R}}
\newcommand{\C}{\mathbb{C}}
\newcommand{\F}{\mathbb{F}}
\newcommand{\A}{\mathbb{A}}
\newcommand{\PP}{\mathbb{P}}
\newcommand{\E}{\mathbb{E}}
\newcommand{\eps}{\varepsilon}
\newcommand{\dd}{\,\mathrm{d}}
\begin{document}
\section*{UnsolvedMath ID: 30003843}
\textbf{Source problem code:} catalogue entry
\subsection{Definitions and mathematical statement}
Use the standard interval-group conventions: let \(I=[0,1]\), and let \(\operatorname{PL}_o(I)\) be the group, under composition, of increasing piecewise-linear homeomorphisms of \(I\) having finitely many linear pieces. Such homeomorphisms fix both endpoints. Let \(F\) denote Thompson's group: the subgroup consisting of maps whose breakpoints lie in \(\Z[1/2]\cap[0,1]\) and whose slopes are integral powers of \(2\). These conventional meanings of \(I\), \(\operatorname{PL}_o(I)\), and \(F\) are not explicitly defined in the supplied record and require confirmation against the cited report.

Write \(G\hookrightarrow H\) when there exists an injective group homomorphism from \(G\) to \(H\). Embeddings here are abstract group embeddings; no compatibility with the given interval actions is imposed. Define biembeddability by
\[
G\sim H\quad\Longleftrightarrow\quad G\hookrightarrow H\ \text{and}\ H\hookrightarrow G.
\]
Let \(\mathcal P\) be the set of biembeddability classes \([G]\) of finitely generated subgroups \(G\leq\operatorname{PL}_o(I)\), ordered by
\[
[G]\leq[H]\quad\Longleftrightarrow\quad G\hookrightarrow H.
\]
Two classes are incomparable if neither is less than or equal to the other. A group is elementary amenable if it belongs to the smallest class containing all finite groups and all abelian groups and closed under taking subgroups, quotients, extensions, and directed unions. Elementary amenability is invariant under biembeddability, so it is meaningful to apply this property to a class in \(\mathcal P\).

Define the following three regions:
\begin{align*}
\mathcal R_{\mathrm{EA}}={}&\{[G]\in\mathcal P:\ G\text{ is elementary amenable},\ G\not\hookrightarrow F,\ F\not\hookrightarrow G\},\\
\mathcal R_{\mathrm{NEA}}={}&\{[G]\in\mathcal P:\ G\text{ is not elementary amenable},\ G\not\hookrightarrow F,\ F\not\hookrightarrow G\},\\
\mathcal R_{<F}={}&\{[G]\in\mathcal P:\ G\hookrightarrow F,\ F\not\hookrightarrow G,\ G\text{ is not elementary amenable}\}.
\end{align*}
The third region can equivalently be represented by finitely generated, non-elementary-amenable subgroups \(G\leq F\) into which \(F\) does not embed: an embedding into \(F\) identifies a representative with its image.

\textbf{Question.} Are all three regions empty; that is, does
\[
\mathcal R_{\mathrm{EA}}=\mathcal R_{\mathrm{NEA}}=\mathcal R_{<F}=\varnothing
\]
hold? Finite generation applies throughout, including to the subgroups representing the third region. The supplied record labels this problem open; that label is a supplied curation claim, not an independently verified mathematical conclusion.
\subsection{Short English statement}
Among finitely generated subgroups of the interval piecewise-linear homeomorphism group, ask whether every group is comparable with Thompson's group \(F\) under abstract embeddability, and whether every non-elementary-amenable subgroup of \(F\) contains an isomorphic copy of \(F\). The interval and group notation use explicitly stated standard conventions that require source confirmation.
\subsection{Sources}
\begin{itemize}
\item Supplied local UnsolvedMath registry record, permanent ID \(30003843\), title \emph{Empty Regions in the Biembeddability Poset}, source problem code \texttt{catalogue entry}, in \emph{Oberwolfach Reports Open Problems}. The record supplies proposed year 2018 and no named proposer.
\item Supplied source citation: \emph{Cohomological and Metric Properties of Groups of Homeomorphisms of R (2019)}, Oberwolfach Reports, DOI \texttt{10.4171/owr/2018/26}, pp.~1579--1633. \url{https://doi.org/10.4171/owr/2018/26}. The supplied background describes Oberwolfach Reports as documenting mathematical workshops and their open-problem discussions.
\item Supplied literature source, labeled \emph{Oberwolfach Report 2019}: \url{https://ems.press/content/serial-article-files/46748}.
\item Supplied literature source, labeled \emph{Oberwolfach Report 2020}: \url{https://oa.tib.eu/renate/server/api/core/bitstreams/13a65f3d-fc78-4e5e-8baf-43e2dfd9102a/content}.
\item Literature assessment dated exactly \texttt{2026-08-21}: the supplied curation states that the emptiness of Regions 2, 4, and 6 remains unresolved and that later work still describes the corresponding Brin--Sapir boundary as unknown. This is supplied dated literature triage derived from the OWR curation, not an independently verified assessment. The time-sensitive status and cited claims require independent verification.
\item The registry labels the statement \texttt{corrected\_verified} and reports a manual semantic, typographic, and LaTeX review. These are supplied review claims; this statement-only draft still requires human mathematical review.
\end{itemize}
\end{document}
